\documentclass[12pt, twoside, italian]{book}
%%%%%%%%%%%%%%%%%%%%%%%%%%%%%%%%%
% PACKAGE IMPORTS
%%%%%%%%%%%%%%%%%%%%%%%%%%%%%%%%%


\usepackage[tmargin=2cm,bmargin=2cm,footskip=.2in, rmargin=1.5in,lmargin=1.5in, margin=0.85in, a4paper]{geometry}
\usepackage{bookmark}
\usepackage{libertinus}
\usepackage{amsmath,amsfonts,amsthm,amssymb,mathtools}
%\usepackage[varbb]{newpxmath}
\usepackage[libertine]{newtxmath}
\usepackage{chemfig}
\usepackage{ragged2e}
\usepackage{xfrac}
\usepackage{mhchem}
\usepackage[italian]{babel}
% \usepackage[Sonny]{fncychap}
\usepackage[makeroom]{cancel}
\usepackage{mathtools}
\usepackage{listing}
\usepackage{enumitem}
\usepackage{hyperref,theoremref}
\hypersetup{
	pdftitle={Espansione in multipolo del campo gravitazionale newtoniano},
	colorlinks=true, linkcolor=doc!90,
	bookmarksnumbered=true,
	bookmarksopen=true,
	pdfauthor={Francesco Sermi}
}
\usepackage[most,many,breakable]{tcolorbox}
\usepackage{xcolor}
\usepackage{varwidth}
\usepackage{varwidth}
\usepackage{etoolbox}
%\usepackage{authblk}
\usepackage{nameref}
\usepackage{multicol,array}
\usepackage{tikz}
\usepackage{tikz-cd}
\usepackage{tikz-3dplot}
\usepackage{tikz}
\usetikzlibrary{3d,decorations.pathmorphing}
\usepackage[ruled,vlined,linesnumbered]{algorithm2e}
\usepackage{import}
\usepackage{xifthen}
\usepackage{pdfpages}
\usepackage{transparent}

\newcommand\mycommfont[1]{\footnotesize\ttfamily\textcolor{blue}{#1}}
\SetCommentSty{mycommfont}
\newcommand{\incfig}[1]{%
    \def\svgwidth{\columnwidth}
    \import{./figures/}{#1.pdf_tex}
}

\usepackage{tikzsymbols}
\usepackage{float}
\usepackage[toc]{appendix}
\renewcommand\qedsymbol{$\square$}
\usepackage{hyperref}
\usepackage{mwe}
\usepackage{fancyhdr}
\usepackage{pgfplots}
\usepackage{bm}
\usepackage{dsfont}
%\usepackage{import}
%\usepackage{xifthen}
%\usepackage{pdfpages}
%\usepackage{transparent}
\definecolor{doc}{RGB}{0,0,0}
\theoremstyle{definition}
\newtheorem{definition}{Definizione}[chapter]

\theoremstyle{definition}
\newtheorem{theorem}{Teorema}[chapter]
\theoremstyle{definition}
\newtheorem{cor}{Corollario}[theorem]
\theoremstyle{definition}
\newtheorem{lemma}{Lemma}[theorem]
\theoremstyle{plain}
\newtheorem{exercise}{Esercizio}[chapter]
\theoremstyle{definition}
\newtheorem{prop}{Proposizione}[chapter]

\theoremstyle{definition}
\newtheorem{example}{Esempio}[chapter]

\theoremstyle{definition}
\newtheorem{axiom}{Assioma}[chapter]

\theoremstyle{remark}
\newtheorem*{remark}{Oss:}

\newcommand*\circled[1]{\tikz[baseline=(char.base)]{
		\node[shape=circle,draw,inner sep=1pt] (char) {#1};}}
\newcommand{\innerprod}[2]{\langle #1, #2 \rangle}
\renewcommand\appendixpagename{Appendici}
\renewcommand\appendixtocname{Appendici}
\setcounter{localmathalphabets}{0}

% Things Lie
\newcommand{\kb}{\mathfrak b}
\newcommand{\kg}{\mathfrak g}
\newcommand{\kh}{\mathfrak h}
\newcommand{\kn}{\mathfrak n}
\newcommand{\ku}{\mathfrak u}
\newcommand{\kz}{\mathfrak z}
\DeclareMathOperator{\Ext}{Ext} % Ext functor
\DeclareMathOperator{\Tor}{Tor} % Tor functor
\newcommand{\gl}{\opname{\mathfrak{gl}}} % frak gl group
\renewcommand{\sl}{\opname{\mathfrak{sl}}} % frak sl group chktex 6

% More script letters etc.
\newcommand{\SA}{\mathcal A}
\newcommand{\SB}{\mathcal B}
\newcommand{\SC}{\mathcal C}
\newcommand{\SF}{\mathcal F}
\newcommand{\SG}{\mathcal G}
\newcommand{\SH}{\mathcal H}
\newcommand{\OO}{\mathcal O}

\newcommand{\SCA}{\mathscr A}
\newcommand{\SCB}{\mathscr B}
\newcommand{\SCC}{\mathscr C}
\newcommand{\SCD}{\mathscr D}
\newcommand{\SCE}{\mathscr E}
\newcommand{\SCF}{\mathscr F}
\newcommand{\SCG}{\mathscr G}
\newcommand{\SCH}{\mathscr H}

% Mathfrak primes
\newcommand{\km}{\mathfrak m}
\newcommand{\kp}{\mathfrak p}
\newcommand{\kq}{\mathfrak q}

% number sets
\newcommand{\RR}[1][]{\ensuremath{\ifstrempty{#1}{\mathbb{R}}{\mathbb{R}^{#1}}}}
\newcommand{\NN}[1][]{\ensuremath{\ifstrempty{#1}{\mathbb{N}}{\mathbb{N}^{#1}}}}
\newcommand{\ZZ}[1][]{\ensuremath{\ifstrempty{#1}{\mathbb{Z}}{\mathbb{Z}^{#1}}}}
\newcommand{\QQ}[1][]{\ensuremath{\ifstrempty{#1}{\mathbb{Q}}{\mathbb{Q}^{#1}}}}
\newcommand{\CC}[1][]{\ensuremath{\ifstrempty{#1}{\mathbb{C}}{\mathbb{C}^{#1}}}}
\newcommand{\PP}[1][]{\ensuremath{\ifstrempty{#1}{\mathbb{P}}{\mathbb{P}^{#1}}}}
\newcommand{\HH}[1][]{\ensuremath{\ifstrempty{#1}{\mathbb{H}}{\mathbb{H}^{#1}}}}
\newcommand{\FF}[1][]{\ensuremath{\ifstrempty{#1}{\mathbb{F}}{\mathbb{F}^{#1}}}}
% expected value
\newcommand{\EE}{\ensuremath{\mathbb{E}}}
\newcommand{\charin}{\text{ char }}
\DeclareMathOperator{\sign}{sign}
\DeclareMathOperator{\Aut}{Aut}
\DeclareMathOperator{\Inn}{Inn}
\DeclareMathOperator{\Syl}{Syl}
\DeclareMathOperator{\Gal}{Gal}
\DeclareMathOperator{\GL}{GL} % General linear group
\DeclareMathOperator{\SL}{SL} % Special linear group

%---------------------------------------
% BlackBoard Math Fonts :-
%---------------------------------------

%Captital Letters
\newcommand{\bbA}{\mathbb{A}}	\newcommand{\bbB}{\mathbb{B}}
\newcommand{\bbC}{\mathbb{C}}	\newcommand{\bbD}{\mathbb{D}}
\newcommand{\bbE}{\mathbb{E}}	\newcommand{\bbF}{\mathbb{F}}
\newcommand{\bbG}{\mathbb{G}}	\newcommand{\bbH}{\mathbb{H}}
\newcommand{\bbI}{\mathbb{I}}	\newcommand{\bbJ}{\mathbb{J}}
\newcommand{\bbK}{\mathbb{K}}	\newcommand{\bbL}{\mathbb{L}}
\newcommand{\bbM}{\mathbb{M}}	\newcommand{\bbN}{\mathbb{N}}
\newcommand{\bbO}{\mathbb{O}}	\newcommand{\bbP}{\mathbb{P}}
\newcommand{\bbQ}{\mathbb{Q}}	\newcommand{\bbR}{\mathbb{R}}
\newcommand{\bbS}{\mathbb{S}}	\newcommand{\bbT}{\mathbb{T}}
\newcommand{\bbU}{\mathbb{U}}	\newcommand{\bbV}{\mathbb{V}}
\newcommand{\bbW}{\mathbb{W}}	\newcommand{\bbX}{\mathbb{X}}
\newcommand{\bbY}{\mathbb{Y}}	\newcommand{\bbZ}{\mathbb{Z}}

%---------------------------------------
% MathCal Fonts :-
%---------------------------------------

%Captital Letters
\newcommand{\mcA}{\mathcal{A}}	\newcommand{\mcB}{\mathcal{B}}
\newcommand{\mcC}{\mathcal{C}}	\newcommand{\mcD}{\mathcal{D}}
\newcommand{\mcE}{\mathcal{E}}	\newcommand{\mcF}{\mathcal{F}}
\newcommand{\mcG}{\mathcal{G}}	\newcommand{\mcH}{\mathcal{H}}
\newcommand{\mcI}{\mathcal{I}}	\newcommand{\mcJ}{\mathcal{J}}
\newcommand{\mcK}{\mathcal{K}}	\newcommand{\mcL}{\mathcal{L}}
\newcommand{\mcM}{\mathcal{M}}	\newcommand{\mcN}{\mathcal{N}}
\newcommand{\mcO}{\mathcal{O}}	\newcommand{\mcP}{\mathcal{P}}
\newcommand{\mcQ}{\mathcal{Q}}	\newcommand{\mcR}{\mathcal{R}}
\newcommand{\mcS}{\mathcal{S}}	\newcommand{\mcT}{\mathcal{T}}
\newcommand{\mcU}{\mathcal{U}}	\newcommand{\mcV}{\mathcal{V}}
\newcommand{\mcW}{\mathcal{W}}	\newcommand{\mcX}{\mathcal{X}}
\newcommand{\mcY}{\mathcal{Y}}	\newcommand{\mcZ}{\mathcal{Z}}


%---------------------------------------
% Bold Math Fonts :-
%---------------------------------------

%Captital Letters
\newcommand{\bmA}{\boldsymbol{A}}	\newcommand{\bmB}{\boldsymbol{B}}
\newcommand{\bmC}{\boldsymbol{C}}	\newcommand{\bmD}{\boldsymbol{D}}
\newcommand{\bmE}{\boldsymbol{E}}	\newcommand{\bmF}{\boldsymbol{F}}
\newcommand{\bmG}{\boldsymbol{G}}	\newcommand{\bmH}{\boldsymbol{H}}
\newcommand{\bmI}{\boldsymbol{I}}	\newcommand{\bmJ}{\boldsymbol{J}}
\newcommand{\bmK}{\boldsymbol{K}}	\newcommand{\bmL}{\boldsymbol{L}}
\newcommand{\bmM}{\boldsymbol{M}}	\newcommand{\bmN}{\boldsymbol{N}}
\newcommand{\bmO}{\boldsymbol{O}}	\newcommand{\bmP}{\boldsymbol{P}}
\newcommand{\bmQ}{\boldsymbol{Q}}	\newcommand{\bmR}{\boldsymbol{R}}
\newcommand{\bmS}{\boldsymbol{S}}	\newcommand{\bmT}{\boldsymbol{T}}
\newcommand{\bmU}{\boldsymbol{U}}	\newcommand{\bmV}{\boldsymbol{V}}
\newcommand{\bmW}{\boldsymbol{W}}	\newcommand{\bmX}{\boldsymbol{X}}
\newcommand{\bmY}{\boldsymbol{Y}}	\newcommand{\bmZ}{\boldsymbol{Z}}
%Small Letters
\newcommand{\bma}{\boldsymbol{a}}	\newcommand{\bmb}{\boldsymbol{b}}
\newcommand{\bmc}{\boldsymbol{c}}	\newcommand{\bmd}{\boldsymbol{d}}
\newcommand{\bme}{\boldsymbol{e}}	\newcommand{\bmf}{\boldsymbol{f}}
\newcommand{\bmg}{\boldsymbol{g}}	\newcommand{\bmh}{\boldsymbol{h}}
\newcommand{\bmi}{\boldsymbol{i}}	\newcommand{\bmj}{\boldsymbol{j}}
\newcommand{\bmk}{\boldsymbol{k}}	\newcommand{\bml}{\boldsymbol{l}}
\newcommand{\bmm}{\boldsymbol{m}}	\newcommand{\bmn}{\boldsymbol{n}}
\newcommand{\bmo}{\boldsymbol{o}}	\newcommand{\bmp}{\boldsymbol{p}}
\newcommand{\bmq}{\boldsymbol{q}}	\newcommand{\bmr}{\boldsymbol{r}}
\newcommand{\bms}{\boldsymbol{s}}	\newcommand{\bmt}{\boldsymbol{t}}
\newcommand{\bmu}{\boldsymbol{u}}	\newcommand{\bmv}{\boldsymbol{v}}
\newcommand{\bmw}{\boldsymbol{w}}	\newcommand{\bmx}{\boldsymbol{x}}
\newcommand{\bmy}{\boldsymbol{y}}	\newcommand{\bmz}{\boldsymbol{z}}

%---------------------------------------
% Scr Math Fonts :-
%---------------------------------------

\newcommand{\sA}{{\mathscr{A}}}   \newcommand{\sB}{{\mathscr{B}}}
\newcommand{\sC}{{\mathscr{C}}}   \newcommand{\sD}{{\mathscr{D}}}
\newcommand{\sE}{{\mathscr{E}}}   \newcommand{\sF}{{\mathscr{F}}}
\newcommand{\sG}{{\mathscr{G}}}   \newcommand{\sH}{{\mathscr{H}}}
\newcommand{\sI}{{\mathscr{I}}}   \newcommand{\sJ}{{\mathscr{J}}}
\newcommand{\sK}{{\mathscr{K}}}   \newcommand{\sL}{{\mathscr{L}}}
\newcommand{\sM}{{\mathscr{M}}}   \newcommand{\sN}{{\mathscr{N}}}
\newcommand{\sO}{{\mathscr{O}}}   \newcommand{\sP}{{\mathscr{P}}}
\newcommand{\sQ}{{\mathscr{Q}}}   \newcommand{\sR}{{\mathscr{R}}}
\newcommand{\sS}{{\mathscr{S}}}   \newcommand{\sT}{{\mathscr{T}}}
\newcommand{\sU}{{\mathscr{U}}}   \newcommand{\sV}{{\mathscr{V}}}
\newcommand{\sW}{{\mathscr{W}}}   \newcommand{\sX}{{\mathscr{X}}}
\newcommand{\sY}{{\mathscr{Y}}}   \newcommand{\sZ}{{\mathscr{Z}}}


%---------------------------------------
% Math Fraktur Font
%---------------------------------------

%Captital Letters
\newcommand{\mfA}{\mathfrak{A}}	\newcommand{\mfB}{\mathfrak{B}}
\newcommand{\mfC}{\mathfrak{C}}	\newcommand{\mfD}{\mathfrak{D}}
\newcommand{\mfE}{\mathfrak{E}}	\newcommand{\mfF}{\mathfrak{F}}
\newcommand{\mfG}{\mathfrak{G}}	\newcommand{\mfH}{\mathfrak{H}}
\newcommand{\mfI}{\mathfrak{I}}	\newcommand{\mfJ}{\mathfrak{J}}
\newcommand{\mfK}{\mathfrak{K}}	\newcommand{\mfL}{\mathfrak{L}}
\newcommand{\mfM}{\mathfrak{M}}	\newcommand{\mfN}{\mathfrak{N}}
\newcommand{\mfO}{\mathfrak{O}}	\newcommand{\mfP}{\mathfrak{P}}
\newcommand{\mfQ}{\mathfrak{Q}}	\newcommand{\mfR}{\mathfrak{R}}
\newcommand{\mfS}{\mathfrak{S}}	\newcommand{\mfT}{\mathfrak{T}}
\newcommand{\mfU}{\mathfrak{U}}	\newcommand{\mfV}{\mathfrak{V}}
\newcommand{\mfW}{\mathfrak{W}}	\newcommand{\mfX}{\mathfrak{X}}
\newcommand{\mfY}{\mathfrak{Y}}	\newcommand{\mfZ}{\mathfrak{Z}}
%Small Letters
\newcommand{\mfa}{\mathfrak{a}}	\newcommand{\mfb}{\mathfrak{b}}
\newcommand{\mfc}{\mathfrak{c}}	\newcommand{\mfd}{\mathfrak{d}}
\newcommand{\mfe}{\mathfrak{e}}	\newcommand{\mff}{\mathfrak{f}}
\newcommand{\mfg}{\mathfrak{g}}	\newcommand{\mfh}{\mathfrak{h}}
\newcommand{\mfi}{\mathfrak{i}}	\newcommand{\mfj}{\mathfrak{j}}
\newcommand{\mfk}{\mathfrak{k}}	\newcommand{\mfl}{\mathfrak{l}}
\newcommand{\mfm}{\mathfrak{m}}	\newcommand{\mfn}{\mathfrak{n}}
\newcommand{\mfo}{\mathfrak{o}}	\newcommand{\mfp}{\mathfrak{p}}
\newcommand{\mfq}{\mathfrak{q}}	\newcommand{\mfr}{\mathfrak{r}}
\newcommand{\mfs}{\mathfrak{s}}	\newcommand{\mft}{\mathfrak{t}}
\newcommand{\mfu}{\mathfrak{u}}	\newcommand{\mfv}{\mathfrak{v}}
\newcommand{\mfw}{\mathfrak{w}}	\newcommand{\mfx}{\mathfrak{x}}
\newcommand{\mfy}{\mathfrak{y}}	\newcommand{\mfz}{\mathfrak{z}}

%From M275 "Topology" at SJSU
\newcommand{\id}{\mathrm{id}}
\newcommand{\taking}[1]{\xrightarrow{#1}}
\newcommand{\inv}{^{-1}}

%From M170 "Introduction to Graph Theory" at SJSU
\DeclareMathOperator{\diam}{diam}
\DeclareMathOperator{\ord}{ord}
\newcommand{\defeq}{\overset{\mathrm{def}}{=}}

%From the USAMO .tex files
\newcommand{\ts}{\textsuperscript}
\newcommand{\dg}{^\circ}
\newcommand{\ii}{\item}

% % From Math 55 and Math 145 at Harvard
% \newenvironment{subproof}[1][Proof]{%
% \begin{proof}[#1] \renewcommand{\qedsymbol}{$\blacksquare$}}%
% {\end{proof}}

\newcommand{\liff}{\leftrightarrow}
\newcommand{\lthen}{\rightarrow}
\newcommand{\opname}{\operatorname}
\newcommand{\surjto}{\twoheadrightarrow}
\newcommand{\injto}{\hookrightarrow}
\newcommand{\On}{\mathrm{On}} % ordinals
\DeclareMathOperator{\img}{im} % Image
\DeclareMathOperator{\Img}{Im} % Image
\DeclareMathOperator{\coker}{coker} % Cokernel
\DeclareMathOperator{\Coker}{Coker} % Cokernel
\DeclareMathOperator{\Ker}{Ker} % Kernel
\DeclareMathOperator{\rank}{rank}
\DeclareMathOperator{\Spec}{Spec} % spectrum
\DeclareMathOperator{\Tr}{Tr} % trace
\DeclareMathOperator{\pr}{pr} % projection
\DeclareMathOperator{\ext}{ext} % extension
\DeclareMathOperator{\pred}{pred} % predecessor
\DeclareMathOperator{\dom}{dom} % domain
\DeclareMathOperator{\ran}{ran} % range
\DeclareMathOperator{\Hom}{Hom} % homomorphism
\DeclareMathOperator{\Mor}{Mor} % morphisms
\DeclareMathOperator{\End}{End} % endomorphism

\newcommand{\eps}{\epsilon}
\newcommand{\veps}{\varepsilon}
\newcommand{\ol}{\overline}
\newcommand{\ul}{\underline}
\newcommand{\wt}{\widetilde}
\newcommand{\wh}{\widehat}
\newcommand{\vocab}[1]{\textbf{\color{blue} #1}}
\providecommand{\half}{\frac{1}{2}}
\newcommand{\dang}{\measuredangle} %% Directed angle
\newcommand{\ray}[1]{\overrightarrow{#1}}
\newcommand{\seg}[1]{\overline{#1}}
\newcommand{\arc}[1]{\wideparen{#1}}
\DeclareMathOperator{\cis}{cis}
\DeclareMathOperator*{\lcm}{lcm}
\DeclareMathOperator*{\argmin}{arg min}
\DeclareMathOperator*{\argmax}{arg max}
\newcommand{\cycsum}{\sum_{\mathrm{cyc}}}
\newcommand{\symsum}{\sum_{\mathrm{sym}}}
\newcommand{\cycprod}{\prod_{\mathrm{cyc}}}
\newcommand{\symprod}{\prod_{\mathrm{sym}}}
\newcommand{\Qed}{\begin{flushright}\qed\end{flushright}}
\newcommand{\parinn}{\setlength{\parindent}{1cm}}
\newcommand{\parinf}{\setlength{\parindent}{0cm}}
% \newcommand{\norm}{\|\cdot\|}
\newcommand{\inorm}{\norm_{\infty}}
\newcommand{\opensets}{\{V_{\alpha}\}_{\alpha\in I}}
\newcommand{\oset}{V_{\alpha}}
\newcommand{\opset}[1]{V_{\alpha_{#1}}}
\newcommand{\lub}{\text{lub}}
\newcommand{\del}[2]{\frac{\partial #1}{\partial #2}}
\newcommand{\Del}[3]{\frac{\partial^{#1} #2}{\partial^{#1} #3}}
\newcommand{\deld}[2]{\dfrac{\partial #1}{\partial #2}}
\newcommand{\Deld}[3]{\dfrac{\partial^{#1} #2}{\partial^{#1} #3}}
\newcommand{\lm}{\lambda}
\newcommand{\uin}{\mathbin{\rotatebox[origin=c]{90}{$\in$}}}
\newcommand{\usubset}{\mathbin{\rotatebox[origin=c]{90}{$\subset$}}}
\newcommand{\lt}{\left}
\newcommand{\rt}{\right}
\newcommand{\bs}[1]{\boldsymbol{#1}}
\newcommand{\exs}{\exists}
\newcommand{\st}{\strut}
\newcommand{\dps}[1]{\displaystyle{#1}}

\newcommand{\sol}{\setlength{\parindent}{0cm}\textbf{\textit{Soluzione:}}\setlength{\parindent}{1cm} }
\newcommand{\solve}[1]{\setlength{\parindent}{0cm}\textbf{\textit{Soluzione: }}\setlength{\parindent}{1cm}#1 \Qed}
\DeclareMathOperator{\Div}{div}
\DeclareMathOperator{\Rot}{rot}
\newcommand{\dcdot}{%
  \mathbin{%
    \nonscript\mspace{-\muexpr\medmuskip*2/3}%
    \cdot
    \nonscript\mspace{-\muexpr\medmuskip*2/3}%
  }%
}

\usepackage{fancyhdr}
\usepackage{titlesec}  % Opzionale, per controllo avanzato sui titoli
\pdfmapline{+rtitillium Titillium-Bold <Titillium-Regular.ttf}
\pagestyle{fancy}
\newcommand{\titillium}{\fontfamily{rtitillium}\selectfont}
% Definizione dell'intestazione per pagine dispari (Odd)
\fancyhead[LO]{\textbf{\thepage}}  % Numero di pagina in grassetto a sinistra
\fancyhead[RO]{\nouppercase\rightmark}         % Capitolo e sezione a destra

% Definizione dell'intestazione per pagine pari (Even)
\fancyhead[LE]{\nouppercase\leftmark}          % Capitolo e sezione a sinistra
\fancyhead[RE]{\textbf{\thepage}}  % Numero di pagina in grassetto a destra

% Cancella il footer
\fancyfoot{}
\setlength{\headheight}{14.49998pt}
\addtolength{\topmargin}{-2.49998pt}
% Personalizzazione della linea orizzontale
\renewcommand{\headrulewidth}{0.4pt}
\renewcommand{\footrulewidth}{0pt}

\pgfplotsset{compat=newest}
\title{Momenti di multipolo del campo gravitazionale}
\author{Francesco Sermi}
\date{\today}
\pgfplotsset{
	compat = newest,
	colormap/violet % se si vuole cambiare il colormap utilizzato per generare i grafici, tuttavia la copertina ha un colormap a
					% parte che va modificato qua sotto. Se si stampa non a colori consiglio blackwhite come colormap
	}
\renewcommand{\appendixpage}{}
\begin{document}
\nocite{poisson, gualtieri, jackson, Holmes2002, Cardoso_2016}
	\begin{titlepage}
    \centering
    
    % Intestazione Università e Dipartimento
    \vspace{0.6cm}
    \vspace{1.2cm}
    
    % Logo istituzionale Unipi
    % Assicurati di avere il file del cherubino/stemma nella stessa cartella
    \includegraphics[width=0.28\textwidth]{Immagini/marchio_unipi.png}\par
    
    \vspace{1.4cm}
    {\large\titillium \large\textbf{DIPARTIMENTO DI FISICA}\par}
    \vspace{0.3cm}
    {\normalsize Corso di Laurea Triennale in Fisica\par}
    

    % Tipologia documento
    \vspace{0.6cm}
    
    % Titolo della Tesi
    {\LARGE \textbf{Momenti di multipolo del campo gravitazionale}\par}
    
    \vfill % Spinge le firme e la data verso il basso mantenendo le proporzioni
    
    % Firme su due colonne
    \begin{minipage}[t]{0.48\textwidth}
        \textbf{Relatore:}\\
        Prof.~Leonardo Gualtieri\\[0.4cm]
        \textbf{Correlatori:} \\ % Rimuovi o commenta se non presente
        Ismaele Bartoluccio \\
	Davide Labrosciano
	\end{minipage}
    \hfill
    \begin{minipage}[t]{0.48\textwidth}
        \begin{flushright}
            \textbf{Candidato:}\\
            Francesco Sermi\\
            Matricola 677094
        \end{flushright}
    \end{minipage}
    
    \vspace{1.8cm}
    
    % Anno Accademico
    \hrule height 0.5pt
    \vspace{0.4cm}
    {\small Anno Accademico 2025/2026\par}
    
\end{titlepage}
	\thispagestyle{empty}
    	\vspace*{8cm}
    	\begin{flushright}
    	\begin{minipage}{0.6\textwidth}
        	\itshape\small
        	\noindent Alla mamma, che mi ha cresciuto senza mai dubitare delle mie capacità. \\
    		\end{minipage}
	\end{flushright}
    	\vfill
    	\newpage
	\chapter*{Long abstract}
In questa tesi si analizza l’espansione in multipoli della gravità come strumento per lo studio del campo gravitazionale dei corpi celesti e delle possibili deviazioni dalla simmetria sferica. Nei modelli di base, pianeti e stelle vengono spesso descritti come sfere perfette e omogenee; nella realtà fisica, la rotazione e le eterogeneità interne generano continue asimmetrie. L’espansione in multipoli consente di superare questa semplificazione, scomponendo il potenziale gravitazionale in una serie di contributi associati a configurazioni geometriche elementari e ben definite. In questo modo è possibile collegare direttamente la geometria del campo gravitazionale alla reale distribuzione della massa all’interno del corpo celeste. \\

Successivamente viene analizzato il legame tra i contributi multipolari e le simmetrie del sistema. Si illustra come per una sfera ideale sopravviva solo il termine di monopolo, pari alla massa totale, mentre nei corpi in rotazione la simmetria assiale selezioni componenti specifiche associate allo schiacciamento ai poli e al rigonfiamento equatoriale. Nei casi asimmetrici più generali intervengono termini di ordine superiore che catturano i dettagli locali. Questa trattazione consente di chiarire il significato fisico dei momenti di multipolo e di utilizzarli come parametri quantitativi per descrivere la struttura del corpo. \\

L’approccio teorico viene successivamente applicato al caso della Terra per mostrare come le irregolarità gravitazionali riflettano la distribuzione non uniforme della massa planetaria. In particolare, viene esaminato il modello di geopotenziale EIGEN-6C4, ricavato dalle osservazioni dei satelliti in orbita, che spinge l’espansione multipolare fino all'ordine 2190. L’analisi di questi contributi permette di mappare le anomalie gravitazionali nelle diverse regioni del globo, collegandole a variazioni di densità nella crosta e nel mantello terrestre, e dimostrando l'efficacia del metodo su un sistema fisico reale. \\

Nella parte conclusiva viene affrontato il problema della deformabilità mareale, considerando la risposta di un corpo celeste all’attrazione gravitazionale di una stella o di un pianeta vicino. Il fenomeno viene descritto nell'ambito newtoniano ed esteso, nel limite di campo debole, alla Relatività Generale. Si dimostra come il campo esterno induca una distorsione della massa quantificabile tramite un coefficiente di deformabilità mareale. \\

Nel complesso, questa tesi illustra l'espansione in multipoli come un ponte tra la teoria del potenziale e l'interpretazione fisica delle proprietà dei corpi celesti. Il percorso sviluppato ne evidenzia l'utilità operativa: dalla comprensione delle simmetrie fondamentali alla geofisica terrestre, fino alla risposta mareale di oggetti compatti, mostrando come ogni ordine multipolare arricchisca progressivamente il quadro descrittivo del sistema.

	\tableofcontents
    	\chapter{Introduzione}
Il campo gravitazionale generato da un corpo racchiude informazioni fondamentali sulla sua struttura interna. Risulta quindi essenziale studiare i contributi derivanti dalle deviazioni dalla simmetria sferica, i quali permettono di rispondere a numerosi interrogativi sulle caratteristiche del pianeta osservato. In questa breve dissertazione descriveremo l'espansione in multipoli nella gravità newtoniana, strumento cardine per lo studio di tali perturbazioni, applicandone i risultati a un sistema soggetto a interazioni mareali. Illustreremo infine come queste stesse configurazioni vengano modellate nel formalismo della Relatività Generale, in cui la gravità perde lo status di forza per divenire espressione della curvatura dello spaziotempo. 
È lecito domandarsi perché affrontare l'argomento in regime newtoniano. Sebbene da un secolo sia noto che non rappresenti la teoria definitiva della gravità, essa costituisce ancora oggi un'approssimazione eccellente in una grande varietà di contesti pratici e astrofisici. Descrive con accuratezza la struttura terrestre e le maree provocate dal Sole e dalla Luna; persino le simulazioni numeriche che tracciano l'evoluzione degli ammassi galattica,la cui formazione è dominata dall'influenza della materia oscura, si basano sulla gravitazione newtoniana. I limiti di questo approccio classico emergono solo quando ci si scontra con la sua incapacità di spiegare la precessione del perielio di Mercurio o la deflessione della luce, fenomeni confermati sperimentalmente e perfettamente coerenti con le previsioni della Relatività Generale. Il criterio fondamentale per stabilire quando sia valido ricorrere all'approssimazione newtoniana è dettato dal fattore di correzione relativistico:
\begin{align}
    \epsilon \sim \frac{GM}{rc^2} \sim \frac{v^2}{c^2},
\end{align}
dove $G$ è la costante gravitazionale universale, $c$ è la velocità della luce, mentre $M$, $r$ e $v$ rappresentano rispettivamente la massa caratteristica, la dimensione (o distanza di separazione) e la velocità del sistema studiato.

\section{Gravitazione newtoniana}

Cominciamo a trarre i nostri primi passi nella teoria newtoniana della gravitazione partendo dalla celeberrima legge di Newton e la legge di gravitazione universale:
\begin{align}
    &m_I \vec{a} = \vec{F} \label{eq:inertial_mass_motion_law}, \\
    &F = - Gm_G M \frac{\vec{r}(t)}{|r(t)|^3}.
\end{align}
Nella prima equazione, $\vec{F}$ è la forza che agisce sul corpo con massa inerziale $m_I$ situato in posizione $\vec{r}(t)$ e $\vec{a} = \frac{d^2\vec{r}}{dt^2}$ la sua accelerazione; mentre nella seconda
equazione la forza $\vec{F}$ è considerata di natura \emph{gravitazionale} ed è orientata nella retta che congiunge i due corpi. Queste equazioni rappresentano le equazioni del moto di un corpo situato in $\vec{r}(t)$ e soggetto alla forza gravitazionale di un'altra massa posta all'origine del nostro sistema di coordinate. La quantità $G$ è detta \emph{costante di gravitazione universale} e, ad oggi, sappiamo che
il suo valore è circa
\begin{align}
    G = (6.6738 \pm 0.0008) \times 10^{-11} \text{m}^2 \text{kg}^{-1} \text{s}^{-2}.
\end{align}
Si noti, innanzitutto, come questa forza dipenda inversamento dal quadrato della distanza fra i due corpi e punti in direzione opposta al versore che punta all'altro corpo. 

Una forma perfettamente equivalente di esprimere queste leggi è ottenuta scrivendo la forza come gradiente di una funzione potenziale, che indicheremo all'interno di questa dissertazione come \emph{potenziale newtoniano}, da cui risulta che
\begin{align}
    U = \frac{GM}{r},
\end{align}
da cui risulta facilmente che
\begin{align}
    (\nabla U)_i &= \frac{\partial U}{\partial x^i} = \frac{\partial U}{\partial r} \frac{\partial r}{\partial x^i} = \nonumber \\
    &= - \frac{GM}{r^2} \frac{1}{2} \frac{1}{\sqrt{\sum\limits_{\substack{j=1 \\ k=1}}^3 x_j x_k \delta_{jk}}} \left( \sum_{\substack{j=1 \\ k=1}}^3 x_k \delta_{ik} + \sum_{\substack{j=1 \\ k=1}}^3 x_j \delta_{ji} \right) = \nonumber \\
    &= - \frac{GM}{r^2} \frac{1}{\cancel{2}} \frac{1}{\sqrt{\sum\limits_{\substack{j=1 \\ k=1}}^3 x_j x_k \delta_{jk}}} \left( \cancel{2} x_i^2 \right) = - \frac{GM}{r^2} \frac{(\vec{r})_i}{|\vec{r}|} = - \frac{GM}{r^3} (\vec{r})_i = (\vec{F})_i.
\end{align}
Riscrivendo in forma vettoriale si ha che
\begin{align}
    \vec{F} = m_G \nabla U,
\end{align}
da cui usando la \eqref{eq:inertial_mass_motion_law}, si ottiene che
\begin{align}
    m_I \vec{a} = m_G \nabla U,
\end{align}
e, nel caso in cui $m_I = m_G$, allora risulta che
\begin{align}
    \label{eq:acc_gravity}
    \vec{a} = \nabla U.
\end{align}

L'assunzione che la massa gravitazione e la massa inerziale sia la stessa è noto come \emph{principio di equivalenza debole}: come conseguenza diretta di questa assunzione, si vede dalla \eqref{eq:acc_gravity} che l'accelerazione a cui è soggetto un corpo
soggetto ad un campo gravitazionale è indipendente dalla massa del corpo stesso. Nel corso della trattazione supporremo che questo principio sia valido\footnote{molti esperimenti, di cui discuterò in appendice, mostrano che la massa inerziale e gravitazionale coincidano a meno di un parte su $1$ triliardo}, e tale principio è anche alla base del ragionamento che
ha condotto Einstein alla Relatività Generale.

Per adesso, come conseguenza dell'assunzione appena esposta e tramite l'introduzione del potenziale newtoniano, le leggi che governano il moto di un corpo soggetto ad un campo gravitazionale sono le seguenti
\begin{align}
    &\vec{a} = \nabla U, \\
    &U = \frac{GM}{r},
\end{align}
tuttavia queste equazioni sono abbastanza limitate per i nostri scopi, in quanto sono valide esclusivamente per un punto materiale situato in $\vec{r}(t)$ soggetto al campo gravitazionale di un'altro corpo situato nell'origine del nostro sistema delle coordinate: 
vorremmo studiare delle situazioni molto più generali, in cui i nostri corpi sono \emph{estesi}; fatti di materia \emph{continua} (solidi, fluidi o gassosi), quindi possono avere una forma e delle dimensioni che sono completamente arbitrari, oltre a poter evolvere in base alla loro
dinamica interna; e, inoltre, vorremmo persino considerare un numero arbitrario di corpi. E' possibile generalizzare queste equazioni, considerando dei sistemi in cui la descrizione della materia è data da un campo di densità di massa $\rho(\vec{x}, t)$, un campo della pressione $p(\vec{x}, t)$ e un campo di velocità
$\vec{v}(\vec{x}, t)$ che, come si vede qua sopra, sono funzioni dello spazio e del tempo e legate fra loro dalle seguenti equazioni
\begin{enumerate}[label=\protect\circled{\arabic*}]
    \item $\frac{\partial \rho}{\partial t} + \vec{\nabla} \cdot (\rho \vec{v}) = 0$ (equazione di continuità);
    \item $\rho \frac{\mathrm{d} \vec{v}}{\mathrm{d}t} = \rho \vec{\nabla} U - \vec{\nabla} p$ (legge di Eulero);
\end{enumerate}
L'equazione che invece lega il comportamento del campo gravitazionale alla distribuzione spaziale della massa è l'equazione di Poisson
\begin{align}
    \nabla^2 U = - 4 \pi G \rho,
\end{align}
dove $\nabla^2 = \sum\limits_{i=1}^3 \frac{\partial^2}{\partial x_i^2}$.

Le equazioni di continuità e di Eulero non saranno ricavate all'interno di questa dissertazione, sebbene discendano, rispettivamente, dalla conservazione del numero di molecole all'interno di ogni elemento di fluido nel suo moto e dalla seconda equazione di Newton scritta per l'elemento infinitesimo di fluido: ci concentreremo, invece, nella derivazione
dell'equazione di Poisson, pertanto nel resto di questa dissertazione, supporremo di lavorare con una data $\rho(\vec{x}, t)$ che abbia fisicamente senso, ovvero soddisfi l'equazione di continuità e di Eulero.

\subsection{Deduzione dell'equazione di Poisson}

Uno dei principi fondamentali su cui si basa la meccanica newtoniana è quello di sovrapposizione delle forze: appare quindi ragionevole suppore che il campo gravitazionale, almeno nell'ottica newtoniana, sia soggetto a tale principio. Assumendo questo, si deve avere che se consideriamo
$N$ corpi di punti materiali che indicheremo con $1, \ldots, N$, allora il campo gravitazionale in un punto $\vec{x}$ sarà dato dalla somma di tutti i singoli potenziali generati da ogni singolo corpo:
\begin{align}
    U(\vec{x}) = \sum_{i=1}^N U_i = G \sum_{i=1}^N \frac{M_i}{|\vec{x} - \vec{x_i}|},
\end{align}
dove con $\vec{x_i}$ indichiamo la posizione dell'$i-$esimo corpo. La generalizzazione alla materia continua di questa equazione è straightforward: supponiamo di avere un sistema descritto da una certa $\rho(t, x')$, ovvero la distribuzione in funzione del tempo e della posizione della massa del sistema studiato:
abbiamo, allora, che il frammento infinitesimo di massa $\mathrm{d} m$ produrrà il campo gravitazionale di un corpo puntiforme, ovvero
\begin{align}
    \mathrm{d}U(\vec{x}, t) = G \frac{\mathrm{d}m}{|\vec{x}-\vec{x}'|},
\end{align} 
dunque, ricordando che la massa infinitesima è data semplicemente dal prodotto del volume infinitesimo $d^3 x'$ con la densità volumica di massa $\rho(\vec{x'}, t)$, possiamo riscrivere l'equazione sopra ottenuta, ricavando che
\begin{align}
    \mathrm{d}U(\vec{x}, t) = G \frac{\rho(\vec{x}', t)}{|\vec{x}-\vec{x}'|} d^3 x',
\end{align}
da cui, integrando, si ottiene che
\begin{align} \label{eq:pot_convolution}
    U(\vec{x}, t) = G \int \frac{\rho(\vec{x'}, t)}{|\vec{x}-\vec{x'}|} d^3 x'.
\end{align}

La \eqref{eq:pot_convolution} mette il potenziale newtoniano $U$ come funzionale di un'arbitraria $\rho$: questo integrale può essere computato, analiticamente o numericamente, data un'arbitraria e generica $\rho$, tuttavia è bene avere presente che, consequenzialmente, la $U$ avrà significato fisico solo se $\rho$ ha significato fisico, ovvero rispetta l'equazione di continuità e la legge di Eulero. La \eqref{eq:pot_convolution}, nonostante ciò, è scritta sotto forma di equazione integrale,
il che non la rende facilmente usufruibile per trovare il potenziale in una data regione dello spazio: possiamo, però, trasformare questa equazione integrale in una differenziale usando la seguente uguaglianza in senso distribuzionale
\begin{align}
        \nabla^2 \frac{1}{|\vec{x}-\vec{x'}|} = -4 \pi \delta^3(\vec{x} - \vec{x'}) \label{eq:green_delta},
\end{align}
che ci accingiamo a dimostrare: prima di farlo, tuttavia, osserviamo che questo risultato ci permette di derivare perfettamente l'equazione di Poisson per la \eqref{eq:pot_convolution}, in quanto
\begin{align}
    \nabla^2 U = G \int \rho(\vec{x}', t) \nabla^2 \frac{1}{|\vec{x}-\vec{x}'|} d^3 x' = G \int \rho(\vec{x}', t) (-4 \pi \delta^3 (\vec{x}-\vec{x}')) d^3 x' = -4 \pi G \rho(\vec{x}, t).
\end{align} 
Si potrebbe anche mostrare che la \eqref{eq:pot_convolution} è soluzione dell'equazione di Poisson: per farlo, osserviamo che l'equazione
\begin{align}
    \nabla^2 U = -4 \pi G \rho,
\end{align}
può essere risolta, usando il formalismo delle funzione di Green, da
\begin{align}
    U(\vec{x}, t) = \int \mathcal{G}(\vec{x}-\vec{x}') \rho(\vec{x}') d^3 x',
\end{align}
con $\mathcal{G}(\vec{x})$ quella distribuzione tale per cui
\begin{align}
    \nabla^2 \mathcal{G}(\vec{x}-\vec{x}') = - 4 \pi G \delta^3(\vec{x}-\vec{x}'),
\end{align}
dunque si vede, confrontando con la \eqref{eq:green_delta}, che
\begin{align}
    \mathcal{G}(\vec{x}-\vec{x}')=\frac{1}{|\vec{x}-\vec{x}'|}.
\end{align}

\subsubsection{Dimostrazione della \eqref{eq:green_delta}}
Ci accingiamo, in questa sezione, di dimostrare la \eqref{eq:green_delta}. Per farlo, è necessario mostrare le due identità di Green:
\begin{lemma}[identità di Green] \hspace{1cm} \\
    Sia $U \in \mathfrak{m}(\mathbb{R}^3)$ e sia $\varphi \in C^2(U), \psi \in C^1(U)$, allora  
    \begin{enumerate}[label=\protect\circled{\arabic*}]
        \item \begin{align}
            \int_U (\psi \nabla^2 \varphi + \nabla \varphi \cdot \nabla \psi)dV = \oint_{\partial U} \psi \frac{\partial \varphi}{\partial n} dS,
        \end{align}
        dove $\hat{n}$ è il versore normale uscente all'elemento di superfice $dS$ e $\partial U$ è il bordo della frontiera;
        \item \begin{align}
            \int_U \left[ \psi \nabla \cdot (\nabla \varphi) - \varphi \nabla \cdot (\nabla \psi)\right]dV = \oint_{\partial U} \left(\psi \frac{\partial \varphi}{\partial n} - \varphi \frac{\partial \psi}{\partial n}\right) dS.
        \end{align} 
    \end{enumerate}
\end{lemma}
\begin{proof} \hspace{1cm} \\
    \circled{1}: si parte ponendo il campo vettoriale $\vec{F} = \psi \nabla \varphi$ e sappiamo che possiamo usare il teorema della divergenza, dunque
    \begin{align}
            \int_U \nabla \cdot \vec{F} d^3 x &= \int_{U} \nabla \cdot \vec{F} d^3 x = \int_U \nabla \cdot (\psi \nabla \varphi) d^3 x = \\
            &= \int_U (\nabla \psi \cdot \nabla \varphi + \psi \nabla^2 \varphi) d^3 x = \int_{\partial U} \vec{F} \cdot \hat{n} dS = \int_{\partial U} \psi \nabla \varphi \cdot \hat{n} dS,
    \end{align}
    ma per definizione di derivata direzionale segue che $\nabla \varphi \cdot \hat{n} = \frac{\partial \varphi}{\partial n}$, dunque
    \begin{align} \label{eq:green_first_id}
        \int_U (\psi \nabla^2 \varphi + \nabla \varphi \cdot \nabla \psi) d^3 x = \oint_{\partial U} \psi (\nabla \varphi \cdot \hat{n}) dS.
    \end{align}
    \circled{2}: si considera l'espressione della prima identità di Green ottenuta per $\vec{F} = \psi \nabla \varphi$ presente in \eqref{eq:green_first_id} e quella per il campo $\vec{F}' = \varphi \nabla \psi$ (ottenuta considerando la \eqref{eq:green_first_id} e scambiando $\varphi \leftrightarrow \psi$), da cui prendendo 
    la differenza fra le due si ottiene, siccome $\nabla \psi \cdot \nabla \varphi = \nabla \varphi \cdot \nabla \psi$, che
    \begin{align}
        \int_U \left[ \psi \nabla \cdot (\nabla \varphi) - \varphi \nabla \cdot (\nabla \psi)\right]dV = \oint_{\partial U} \left(\psi \frac{\partial \varphi}{\partial n} - \varphi \frac{\partial \psi}{\partial n}\right) dS.
    \end{align} 
\end{proof}
\begin{proof}[Dimostrazione della \eqref{eq:green_delta}]
    Si osserva che $\forall \vec{x}, \vec{x'} \in \mathbb{R}^3, \vec{x} \neq \vec{x'} \implies |\vec{x} - \vec{x'}| \neq 0$, il che implica naturalmente che il gradiente è ben definito per quei punti, dunque
    \begin{align}
        \nabla \frac{1}{|\vec{x}-\vec{x'}|} = \left( -\frac{x_i - x'_i}{|\vec{x}-\vec{x'}|^3} \right)_i,
    \end{align}
    da cui si ricava facilmente, sempre perché $|\vec{x}-\vec{x'}| \neq 0$, che
    \begin{align}
        \nabla \cdot \left( \nabla \frac{1}{|\vec{x}-\vec{x'}|} \right) = \nabla^2 \frac{1}{|\vec{x}-\vec{x'}|},
    \end{align}
    e osserviamo che $\partial_i (\nabla \frac{1}{|\vec{x}-\vec{x'}|})_i = \partial_i \frac{x_i}{|\vec{x}-\vec{x'}|^3} = -\frac{1}{|\vec{x} - \vec{x'}|^3} + \frac{3(x_i-x'_i)^2}{|\vec{x}-\vec{x'}|^5}$, da cui è chiaro che
    \begin{align}
        \nabla^2 \frac{1}{|\vec{x}-\vec{x'}|} &= \sum_{i=1}^3 \partial_i^2 \frac{1}{|\vec{x}-\vec{x'}|} = \sum_{i=1}^3 \left[ \frac{3(x_i - x_i')^2}{|\vec{x}-\vec{x'}|^5} - \frac{1}{|\vec{x}-\vec{x'}|^3} \right] = \nonumber \\
        &= \frac{3}{|\vec{x}-\vec{x'}|^5} \sum_{i=1}^3 [(x_i - x_i')^2] - \frac{3}{|\vec{x}-\vec{x'}|^3} = \frac{3}{|\vec{x} - \vec{x'}|^{\cancel{5} \, 3}} \cancel{|\vec{x}-\vec{x'}|^2} - \frac{3}{|\vec{x} - \vec{x'}|^3} = 0.
    \end{align}
    Osserviamo, tuttavia, che il laplaciano di questa funzione è ben definito se e solo se $|\vec{x} - \vec{x'}| = 0 \iff \vec{x} = \vec{x'}$, come atteso dal breve commento a inizio dimostrazione, pertanto non è possibile concludere la dimostrazione usando la continuità dell'operazione di derivata nello spazio delle distribuzioni. Si procede, quindi, in maniera differente, osservando che possiamo approssimare la funzione $\nabla^2 \frac{1}{|\vec{x} - \vec{x'}|}$ con la funzione così definita
    \begin{align}
            u_{\varepsilon, R} = \begin{cases} 
                \frac{1}{|\vec{x} - \vec{x}'|} & \varepsilon \leq |\vec{x}-\vec{x'}| \leq R, \\
                0 & \text{ altrimenti },
            \end{cases},
    \end{align}
    da cui si ha che usando la seconda identità di Green
    \begin{align}
        \int_{\mathbb{R}^3} (\nabla^2 u_{\varepsilon, R}) \phi(x) d^3 x = \int_{V_{\varepsilon, R}} (\nabla^2 \phi(\vec{x})) - \oint_{\partial V_{\varepsilon, R}} \phi(\vec{x}) \frac{\partial}{\partial n} \left( \frac{1}{|\vec{x}-\vec{x'}|} \right) - \frac{\partial \phi(\vec{x})}{\partial n} \frac{1}{|\vec{x}-\vec{x'}|}dS,
    \end{align}
    tuttavia osserviamo che il bordo dell'insieme $V_{\varepsilon, R}$ è costituito dalla sfera $S_\varepsilon$ di raggio $\varepsilon$ e dalla sfera $S_R$ di raggio $R$ centrate in $|\vec{x}-\vec{x'}|$ e, per il buon andamento delle funzioni di test, si conclude che, nel limite di $R \to +\infty$, l'integrale sulla sfera $S_R$ deve necessariamente fare zero, mentre nel volume $V_{\varepsilon, R}$ si ha che $\nabla^2 \frac{1}{|\vec{x} - \vec{x'}|}$ è ben definito, dunque $\nabla^2 \frac{1}{|\vec{x} - \vec{x'}|} = 0$. Si ha quindi che
    \begin{align}
        \lim_{\substack{R \to +\infty, \\ \varepsilon \to 0}} \int_{\mathbb{R}^3} (\nabla^2 u_{\varepsilon, R}) \phi(\vec{x}) \, d^3x = \lim_{\substack{R \to +\infty \\ \varepsilon \to 0}} \oint_{\partial S_\varepsilon} \left[ \phi(\vec{x}) \frac{\partial}{\partial n} \left( \frac{1}{|\vec{x}-\vec{x'}|} \right) - \frac{\partial \phi(\vec{x})}{\partial n} \frac{1}{|\vec{x}-\vec{x'}|} \right] dS.
    \end{align}
    Definiamo il vettore distanza $\vec{r} = \vec{x} - \vec{x'}$ e il suo modulo $r = |\vec{r}|$. Sulla superficie della sfera interna $\partial S_\varepsilon$, si ha $r = \varepsilon$.

È fondamentale prestare attenzione al versore normale $\hat{n}$: affinché valga la seconda identità di Green, $\hat{n}$ deve essere uscente dal volume $V_{\varepsilon, R}$. Sulla frontiera interna $\partial S_\varepsilon$, tale normale punta dunque verso l'origine delle coordinate sferiche (cioè verso $\vec{x'}$). Si ha quindi $\hat{n} = -\hat{r}$, dove $\hat{r} = \frac{\vec{r}}{r}$ è il versore radiale.

Calcoliamo la derivata direzionale normale della funzione $\frac{1}{|\vec{x}-\vec{x'}|}$:
\begin{align}
    \frac{\partial}{\partial n} \left( \frac{1}{r} \right) = \nabla \left( \frac{1}{r} \right) \cdot \hat{n} = \left( -\frac{1}{r^2} \hat{r} \right) \cdot (-\hat{r}) = \frac{1}{r^2}.
\end{align}
Valutando questa espressione su $\partial S_\varepsilon$ (dove $r = \varepsilon$), essa vale semplicemente $\frac{1}{\varepsilon^2}$. Il secondo termine dell'integranda contiene invece la funzione non derivata, che su $\partial S_\varepsilon$ è $\frac{1}{\varepsilon}$.

Ricordando che l'elemento di superficie in coordinate sferiche è $dS = \varepsilon^2 d\Omega$ (con $d\Omega = \sin\theta \, d\theta \, d\phi$ angolo solido infinitesimo), possiamo riscrivere l'integrale di superficie in questo modo:
\begin{align}
    \oint_{\partial S_\varepsilon} \left[ \phi(\vec{x}) \left( \frac{1}{\varepsilon^2} \right) - \frac{\partial \phi(\vec{x})}{\partial n} \left( \frac{1}{\varepsilon} \right) \right] \varepsilon^2 d\Omega = \oint_{\partial S_\varepsilon} \phi(\vec{x}) \, d\Omega - \varepsilon \oint_{\partial S_\varepsilon} \frac{\partial \phi(\vec{x})}{\partial n} \, d\Omega.
\end{align}

Prendendo ora il limite per $\varepsilon \to 0$, osserviamo che, essendo $\phi$ una funzione di test ($\phi \in \mathcal{S}$), le sue derivate sono ovunque limitate. Pertanto, il secondo integrale è moltiplicato per $\varepsilon$ e tende banalmente a zero:
\begin{align}
    \lim_{\varepsilon \to 0} \varepsilon \oint_{\partial S_\varepsilon} \frac{\partial \phi(\vec{x})}{\partial n} \, d\Omega = 0.
\end{align}

Per il primo termine, grazie alla continuità di $\phi$, al tendere di $\varepsilon \to 0$ la sfera collassa nel punto $\vec{x'}$, e $\phi(\vec{x}) \to \phi(\vec{x'})$. Poiché l'integrale dell'angolo solido su tutta la sfera restituisce $4\pi$, otteniamo:
\begin{align}
    \lim_{\varepsilon \to 0} \oint_{\partial S_\varepsilon} \phi(\vec{x}) \, d\Omega = \phi(\vec{x'}) \oint d\Omega = 4\pi \phi(\vec{x'}).
\end{align}

Infine, ricordando la definizione di derivata in senso distribuzionale e che $\nabla^2 (1/|\vec{x}-\vec{x'}|) = 0$ strettamente all'interno di $V_{\varepsilon, R}$, l'azione del Laplaciano sulla funzione di test si ricava dall'identità di Green portando i termini di superficie a membro opposto (da cui emerge un segno meno):
\begin{align}
    \left\langle \nabla^2 \frac{1}{|\vec{x}-\vec{x'}|}, \phi \right\rangle &= \int_{\mathbb{R}^3} \frac{1}{|\vec{x}-\vec{x'}|} \nabla^2 \phi(\vec{x}) \, d^3x = \nonumber \\
    &= - \lim_{\varepsilon \to 0} \oint_{\partial S_\varepsilon} \left[ \phi(\vec{x}) \frac{\partial}{\partial n} \left( \frac{1}{|\vec{x}-\vec{x'}|} \right) - \frac{\partial \phi(\vec{x})}{\partial n} \frac{1}{|\vec{x}-\vec{x'}|} \right] dS = \nonumber \\
    &= -4\pi \phi(\vec{x'}).
\end{align}
Poiché l'azione della delta di Dirac è definita da $\langle \delta^3(\vec{x} - \vec{x'}), \phi \rangle = \phi(\vec{x'})$, ne consegue la tesi del teorema:
\begin{align}
    \nabla^2 \frac{1}{|\vec{x}-\vec{x'}|} = -4\pi \delta^3(\vec{x} - \vec{x'}).
\end{align}

\end{proof}
	\chapter{Corpi non sferici}
Nel capitolo precedente abbiamo visto che un sistema caratterizzato da una densità $\rho(\vec{x}, t)$ soddisfa l'equazione di Poisson
\begin{align}\nabla^2 U(\vec{x}, t) = -4 \pi G \rho(\vec{x}, t)
\end{align}
che in questo capitolo ci prefiggiamo di risolvere nel caso in cui vi siano deviazioni dalla simmetria sferica. È istruttivo studiare inizialmente il caso con simmetria sferica perfetta per poi passare al caso più generale.
\section{Sistema a simmetria sferica}
Riscriviamo il laplaciano in coordinate sferiche e lo inseriamo nell'equazione di Poisson:
\begin{align}\left[ \frac{1}{r^2} \frac{\partial}{\partial r}\left(r^2 \frac{\partial}{\partial r}\right) + \frac{1}{r^2 \sin{\theta}} \frac{\partial}{\partial \theta}\left(\sin{\theta} \frac{\partial}{\partial \theta}\right) + \frac{1}{r^2 \sin^2{\theta}} \frac{\partial^2}{\partial \phi^2} \right] U = - 4 \pi G \rho.\end{align}
Se il sistema gode di simmetria sferica, il potenziale dipende esclusivamente dalla coordinata radiale $r$; pertanto l'equazione si riduce a
\begin{align}
    \frac{1}{r^2} \frac{\partial}{\partial r} \left( r^2 \frac{\partial U}{\partial r} \right) = -4 \pi G \rho(r, t),
\end{align}
da cui si ottiene
\begin{align}
    \frac{\partial U}{\partial r} = -\frac{4 \pi G}{r^2} \int_0^r \rho(r', t) (r')^2 dr' + \frac{C(t)}{r^2}.
\end{align}
La costante di integrazione $C(t)$ va posta a zero. Fisicamente, infatti, $\frac{\partial U}{\partial r}$ rappresenta il campo gravitazionale radiale. Nell'origine ($r=0$) la direzione radiale non è definita: una forza netta non nulla in quel punto punterebbe in una direzione preferenziale, rompendo la simmetria sferica del sistema. Inoltre, se $C(t)$ fosse diversa da zero, il termine divergerebbe. In definitiva si ha:
\begin{align}
    \frac{\partial U}{\partial r} = - \frac{4 \pi G m(r, t)}{r^2},
\end{align}
dove abbiamo definito la massa contenuta entro il raggio $r$:\begin{align}m(r, t) = \int_0^r \rho(r', t) (r')^2 dr'.\end{align}
Assumendo che il corpo sia confinato nello spazio, esiste un raggio $R$ tale per cui $\rho(r) = 0$ per ogni $r > R$. Pertanto, la massa totale
\begin{align}
    M \equiv m(R, t) = \int_0^R \rho(r', t) (r')^2 dr'
\end{align}
è una costante nel tempo. Lo si dimostra facilmente:
\begin{align}
    \frac{\mathrm{d} M}{\mathrm{d}t} = \int_0^R \frac{\partial \rho(r', t)}{\partial t} (r')^2 dr' = \int_0^R -\frac{1}{(r')^2} \nabla_r (\rho(r', t) v_r)(r')^2 dr' = 0,
\end{align}
dove abbiamo usato l'equazione di continuità in coordinate sferiche e il fatto che la densità e la velocità radiale $v_r$ siano nulle sul bordo del volume di integrazione. 
Abbiamo quindi mostrato che la massa totale del sistema si conserva. Inoltre, poiché per $r > R$ la densità si annulla, il campo gravitazionale all'esterno del corpo è dato da:
\begin{align}
    \frac{\partial U}{\partial r} = - \frac{4 \pi G M}{r^2} \implies U(r) = \frac{4 \pi G M}{r}.
\end{align}
Il potenziale all'esterno coincide con quello generato da una massa puntiforme $M$ posta nel centro del sistema: integrando l'equazione di Poisson abbiamo ottenuto una dimostrazione molto più diretta del teorema del guscio sferico di Newton.

\section{Deviazioni dalla simmetria sferica}
La simmetria sferica è una semplificazione utile, ma è raro che i corpi celesti siano perfettamente simmetrici. La rotazione attorno a un asse, l'interazione con altri corpi e le tensioni interne (come gli stress crostali delle stelle di neutroni) causano deviazioni da questa configurazione ideale. Quando tali deviazioni sono piccole, si procede tramite un'espansione in multipoli. Poiché il contributo dei momenti di ordine superiore tende a decrescere rapidamente, un troncamento dello sviluppo ai primi termini offre una buona descrizione del potenziale gravitazionale per la maggior parte dei problemi di interesse.
Nel caso della Terra, la misurazione dei momenti di multipolo fino a $l = 360$ (che quantificano le variazioni del campo gravitazionale su scale di centinaia di chilometri) ha permesso di studiare fenomeni come la deriva dei continenti, il vulcanismo e la struttura interna del pianeta. Questo approccio analitico, tuttavia, è limitato a sistemi approssimativamente sferici: le galassie a spirale, ad esempio, si discostano marcatamente da tale geometria, obbligando a ricorrere a metodi analitici differenti o a simulazioni numeriche per ricavarne il potenziale.
\subsection{Decomposizione in armoniche sferiche}
Ricordiamo che le armoniche sferiche $Y_{lm}$ sono autofunzioni dell'operatore differenziale angolare:
\begin{align}
    \left[ \frac{1}{\sin{\theta}} \frac{\partial}{\partial \theta} \left( \sin{\theta} \frac{\partial}{\partial \theta} \right) + \frac{1}{\sin^2{\theta}} \frac{\partial^2}{\partial \phi^2} \right] Y_{lm} = -l(l+1) Y_{lm},
\end{align}con $l$ intero non negativo e $m \in \{-l, \dots, l\}$ e riconosciamo, nel termine a sinistra, la parte angolare del laplaciano.
Dalla teoria di Sturm-Liouville sappiamo che queste funzioni costituiscono una base ortogonale e completa per lo spazio $L^2(S^2)$, ossia lo spazio delle funzioni a quadrato integrabile definite sulla sfera. Pertanto, vale la relazione di ortonormalità:
\begin{align}
    \int_{S^2} Y_{lm}^(\theta, \phi) Y_{l'm'}(\theta, \phi) \mathrm{d}\Omega = \delta_{ll'} \delta_{mm'},
\end{align}
e per ogni funzione $f \in L^2(S^2)$ si ha:
\begin{align}
f(\theta, \phi) = \sum_{l=0}^{\infty} \sum_{m=-l}^l f_{lm} Y_{lm}(\theta, \phi),
\end{align}
con i coefficienti dati da:
\begin{align}
f_{lm} = \int_{S^2} Y_{lm}(\theta, \phi) f(\theta, \phi) \mathrm{d}\Omega.
\end{align}
Possiamo quindi sviluppare sia la densità che il potenziale in serie di armoniche sferiche, ottenendo
\begin{align}
    &\rho(r, \theta, \phi, t) = \sum_{l=0}^{+\infty} \sum_{m=-l}^l \rho_{lm}(r, t) Y_{lm}(\theta, \phi), & &U(r, \theta, \phi, t) = \sum_{l=0}^{+\infty} \sum_{m=-l}^l U_{lm}(r, t) Y_{lm}(\theta, \phi),
\end{align}
dove i coefficienti dipendono da $r$ e $t$, poiché la proiezione viene effettuata su superfici sferiche a raggio e tempo fissati, e sono dati da
\begin{align}
    &\rho_{lm}(r, t) = \int_{S^2} Y_{lm}(\theta, \phi) \rho(r, \theta, \phi, t) \mathrm{d}\Omega, &
    &U_{lm}(r, t) = \int_{S^2} Y_{lm}(\theta, \phi) U(r, \theta, \phi, t) \mathrm{d}\Omega.
\end{align}

Inserendo questi sviluppi nell'equazione di Poisson si ottiene:
\begin{align}
    &\left[ \frac{1}{r^2} \frac{\partial}{\partial r}\left(r^2 \frac{\partial}{\partial r}\right) + \frac{1}{r^2 \sin{\theta}} \frac{\partial}{\partial \theta}\left(\sin{\theta} \frac{\partial}{\partial \theta}\right) + \frac{1}{r^2 \sin^2{\theta}} \frac{\partial^2}{\partial \phi^2} \right] \sum_{l=0}^{+\infty} \sum_{m=-l}^l U_{lm} Y_{lm} = \nonumber \\
    &= -4 \pi G \sum_{l=0}^{+\infty} \sum_{m=-l}^l \rho_{lm} Y_{lm}.\end{align}
Poiché le $Y_{lm}$ sono linearmente indipendenti, l'uguaglianza deve valere termine a termine per ogni coppia $(l, m)$, portando a:
\begin{align}
    \left[ \frac{\partial}{\partial r} \left( r^2 \frac{\partial}{\partial r} \right) - l(l+1) \right] U_{lm}(r, t) = - 4 \pi G r^2 \rho_{lm}(r, t).
\end{align}
In questo modo abbiamo ricondotto un'equazione alle derivate parziali a un sistema infinito di equazioni differenziali ordinarie disaccoppiate. Per risolvere queste equazioni utilizziamo il metodo delle funzioni di Green, definendo l'operatore differenziale
\begin{align}
    \mathcal{L} = \frac{\partial}{\partial r} r^2 \frac{\partial}{\partial r} - l(l+1),
\end{align}cerchiamo una funzione $g_l(r, r')$ che soddisfi
\begin{align}
    \mathcal{L} g_l(r, r') = - 4 \pi \delta(r-r').
\end{align}
Una volta determinata $g_l$, ogni coefficiente $U_{lm}$ potrà essere calcolato come
\begin{align} \label{eq:u_lm_green}
    U_{lm}(r, t) = G \int_0^\infty g_l(r, r') \rho_{lm}(r', t) (r')^2 dr'.
\end{align}
L'operatore $\mathcal{L}$ e, di conseguenza, la funzione di Green dipendono dall'indice $l$ ma non da $m$. Per $r \neq r'$, l'equazione per $g_l$ si riduce a quella omogenea:
\begin{align}
    \mathcal{L} g_l(r, r') = 0.
\end{align}
Cercando soluzioni della forma $r^\gamma$, si ha
\begin{align}
    \left[ \frac{\partial}{\partial r} r^2 \frac{\partial}{\partial r} - l(l+1) \right] r^\gamma = 0 \implies \gamma(\gamma+1) - l(l+1) = 0.
\end{align}
La combinazione lineare delle due radici indipendenti fornisce:\begin{align}g_l(r, r') = A(r')r^l + B(r')r^{-(l+1)}.\end{align}I coefficienti $A$ e $B$ dipendono da $r'$ poiché tale variabile gioca il ruolo di parametro per l'equazione differenziale in $r$. Affinché il potenziale sia fisicamente accettabile, non deve presentare divergenze spurie. Per $r < r'$, richiediamo che la soluzione sia regolare nell'origine ($r \to 0$); pertanto imponiamo $B = 0$:\begin{align}g_<(r, r') = A(r')r^l.\end{align}Analogamente, per $r > r'$, richiediamo che il potenziale si annulli all'infinito ($r \to \infty$); pertanto imponiamo $A = 0$:
\begin{align}
    g_>(r, r') = B(r')r^{-(l+1)}.
\end{align}
Combinando i risultati mediante il gradino di Heaviside, si ottiene:
\begin{align}
    g_l(r, r') = A(r')r^l \theta(r'-r) + B(r')r^{-(l+1)}\theta(r-r').
\end{align}
Imponendo infine la continuità della funzione di Green in $r = r'$, si ricava la condizione
\begin{align}
    A(r')(r')^l = B(r')(r')^{-(l+1)},
\end{align}
dunque possiamo riscrivere nella seguente forma
\begin{align}
    g_l(r, r') = C(r') \left[ \frac{r^l}{(r')^l}\theta(r'-r) + \frac{r^{-(l+1)}}{(r')^{-(l+1)}}\theta(r-r') \right].
\end{align}
ed integrando l'equazione differenziale fra $[r'-\varepsilon, r'+\varepsilon]$ con $\varepsilon \to 0$ si ottiene che
\begin{align}
    \int_{r'-\varepsilon}^{r'+\varepsilon} \left[ \frac{\partial}{\partial r} r^2 \frac{\partial}{\partial r} \right] g_l(r, r') = - 4 \pi.
\end{align} 
Si noti che il termine proporzionale all'integrale di $g_l(r, r')$ si annulla per continuità nel limite di $\varepsilon \to 0$: rimane dunque
\begin{align}
    \int_{r'-\varepsilon}^{r'+\varepsilon} \left[ \frac{\partial}{\partial r} \left( r^2 \frac{\partial g_l(r, r')}{\partial r} \right) \right] = - 4 \pi,
\end{align}
in cui si ottiene che
\begin{align}
    C(r')[(-l+1)(r') - (l+2)(r')] = -4 \pi \implies C(r') = \frac{4 \pi}{2l+1} \frac{1}{r'}.
\end{align}
La funzione di Green dell'operatore $\mathcal{L}$ risulta, dunque, essere
\begin{align}
    g_l(r, r') = \frac{4 \pi}{2l+1} \left[ \frac{r^l}{(r')^{l+1}} \theta(r' - r) + \frac{(r')^{l}}{r^{l+1}} \theta(r-r') \right],
\end{align}
che può essere riscritta in forma più compatta introducendo $r_> = \max\{r, r'\}$ e $r_< = \min\{r, r'\}$, ottenendo
\begin{align}
    g_l(r, r') = \frac{4 \pi}{2l+1} \frac{r_<^l}{r_>^{l+1}}.
\end{align}
Sostituendo nella \eqref{eq:u_lm_green} si ottiene che
\begin{align}
    U_{lm}(r, t) &= \frac{4 \pi G}{2l+1} \int_{0}^{+\infty} \frac{r^l_<}{r_>^{l+1}} \rho_{lm}(r', t) \mathrm{d}r' = \nonumber \\
    &= \frac{4 \pi G}{2l+1} \left[ r^l \int_r^{+\infty} \frac{\rho_{lm}(r', t)}{(r')^{l+1}} (r')^2 dr' + \frac{1}{r^{l+1}} \int_0^r (r')^l \rho_{lm}(r', t)(r')^2 \mathrm{d}r' \right].
\end{align}


	\chapter{Deformazioni mareali}

Uno dei fenomeni di maggiore interesse che viene studiato tramite l'espansione in multipoli è la deformazione mareale di un corpo soggetto ad un campo di marea esterno (come quello fra la Luna ela Terra).

\section{Deformazioni mareali in Gravità Newtoniana}

Per un oggetto stazionario, isolato sappiamo che la soluzione è data dalla \eqref{eq:pot_convolution} e lontano dalla sorgente possiamo utilizzare l'espansione multipolare in termini dei tensori STF, per cui risulta che
\begin{align}
	\frac{1}{|\vec{x} - \vec{x}'|} = \frac{1}{r} + \frac{n_i x^{'i}}{r^2} + \frac{3n_i n_j - \delta_{ij}}{2r^3} x^{' i} x^{' j} + o \left( \frac{1}{r^3} \right),
\end{align}
da cui il potenziale risulta essere pari a
\begin{align}
	\Phi(\vec{x}) &= - \frac{1}{r} \int_V \rho(\vec{x}') \mathrm{d}^3 x' - \frac{n_i}{r^2} \int_V \rho(\vec{x}') x^{' i} \mathrm{d}^3 x' - \frac{3}{2} \frac{n_{\langle i} n_{j \rangle}}{r^3} \int \rho(\vec{x}') x^{' i} x^{' j} \mathrm{d}^3 x' + o \left( \frac{1}{r^3} \right) = \\
		      &= - \frac{M}{r} - \frac{3}{2} \frac{Q_{ij} n^{\langle i} n^{j \rangle}}{r^3} + o \left( \frac{1}{r^4} \right),
\end{align}
dove abbiamo riconosciuto che il primo integrale è semplicemente la massa totale del corpo, il secondo integrale, invece, è nullo nel sistema del centro di massa e $Q_{ij}$ è semplicemente il momento di quadrupolo del corpo. Consideriamo adesso un campo stazionario esterno $\Phi_{\text{ext}}$ che espandiamo in serie di Taylor attorno al centro di massa del corpo, posizionato in $\vec{x} = \vec{0}$
\begin{align}
	\Phi_{\text{ext}}(\vec{x}) = \text{const.} + \frac{\partial \Phi_{\text{ext}}}{\partial x^i}\Bigg|_{\vec{x}=\vec{0}} + \frac{1}{2} E_{ij} x^i x^j + o(r^2),
\end{align}
dove $E_{ij}$ è la matrice hessiana del potenziale calcolata in $\vec{x}=\vec{0}$, esplicitamente
\begin{align}
	E_{ij} = \frac{\partial^2 \Phi_{\text{ext}}}{\partial x^i \partial x^j} \Bigg|_{\vec{x}=\vec{0}},
\end{align}
che chiamiamo \emph{tensore tidiale} del campo gravitazionale: questo è il contributo che determina le maree, per esempio, indotte dalla Luna sulla Terra. Se la forza esercitata al centro di massa del corpo sparisce, ovvero $\frac{\partial \Phi_{\text{ext}}}{\partial x^i} = 0$ per $\vec{x}=\vec{0}$, allora il campo esterno è detto campo tidiale, in quanto l'ordine dominate è dato dal tensore tidiale $E_{ij}$.

Si osservi che le due espansioni che abbiamo fatto sono valide esclusivamente se $R \ll r \ll \mathcal{R}$, dove $R$ è la lunghezza caratteristica del corpo e $\mathcal{R}$ è la lunghezza caratteristica del corpo che genera il campo gravitazionale tidiale: chiameremo questa regione \emph{buffer region}, in cui la sorgente del campo gravitazionale mareale è esterna a tale regione, pertanto abbiamo che (tenendo solo il secondo ordine che dà contributo non nullo)
\begin{align}
	\nabla^2 \Phi_{\text{ext}} &= \delta_{j}^i \partial_{i} \partial_j(\frac{1}{2} E_{ab} x^a x^b) = \delta_j^i \partial_i ( \frac12 E_{ab} \delta_j^a x^b + \frac12 E_{ab} x^a \delta_j^b ) = \nonumber \\
	&= \delta_j^i \partial_i ( \frac12 E_{jb} x^b + \frac12 E_{aj} x^a) = \delta_j^i \partial_i (E_{aj} x^a) = \delta_j^i E_{aj} \delta_i^a = E_{ii} = 0,
\end{align}
ovvero si ha che $E_{ij}$ è un tensore simmetrico (per il teorema di Schwartz) e senza traccia. Per linearità dell'equazione di Poisson, abbiamo che il potenziale gravitazionale nella regione di accrescimento è dato semplicemente dalla somma dei due potenziali, ovvero
\begin{align} \label{eq:expansion_buffer}
	\Phi_{\text{tot}}(\vec{x}) = - \frac{M}{r} - \frac32 \frac{Q_{ij} n^{\langle i} n^{j \rangle}}{r^3} + o \left( \frac{R^3}{r^3} \right) + \frac12 E_{ij} r^2 n_{\langle i} n^{j \rangle} + o \left( \frac{r^3}{\mathcal{R}^3} \right), 
\end{align}
dove si osservi che $x^{i} x^j = \frac{x^i x^j}{r^2} r^2 = n^i n^j r^2$, ma poiché $E_{ij}$ è tracefree, allora si ha che $E_{ij} n^i n^j r^2 = E_{ij} n^{\langle i} n^{j \rangle}$.

Questa espansione si rompe vicino al corpo (in cui $r \sim R$), dove i termini di alto ordine nell'espansione in multipoli in $\frac{R}{r}$ non sono trascurabili; inoltre lontano dal corpo, in cui $r \sim \mathcal{R}$, il campo esterno non può essere descritto solo in termini del tensore tidiale e gli altri termini dell'espansioni di alto grado sono anch'essi non trascurabili. Per trattare gli effetti che il campo mareale esercita sul corpo si rivela necessario andare a risolvere l'equazione non solo nel vuoto, ma anche all'interno del corpo, includendo il campo esterno tramite delle accurate condizioni al contorno: facendo questo, si ricava che il corpo acquista un momento di quadrupolo che è proporzionale al campo mareale; esplicitamente
\begin{align}
	Q_{ij} = - \lambda E_{ij},
\end{align}
dove la costante di proporzionalità $\lambda$ è noto in letteratura come la deformabilità mareale quadrupolare, che dipende dalla composizione interna del corpo. Questo problema fu studiato all'inizio del ventesimo secolo da A. Love, il quale notò che la deformabilità mareale quadrupolare ha le dimensioni $[L]^5$ ed è proporzionale a $R^5$ di una stella e definì una quantità adimensionale, nota come \emph{numero di Love} di una stella:
\begin{align}
	\kappa_2 = \frac32 \frac{\lambda}{R^5},
\end{align}
che per le stelle di neutroni assume un valore compreso fra $10^{-2}$ e $10^{-1}$, a seconda dell'equazione di stato di essa e sulla sua compattezza.

\section{Deformazioni mareali in Relatività Generale}

Consideriamo il caso, adesso, di un oggetto descritto dalla Relatività Generale (ovvero con un campo gravitazionale ``forte'', ovvero $\frac{GM}{Rc^2} \gtrsim 1$ con $R$ scala caratteristica delle lunghezze del corpo): lontano da esso, ovvero per $r \gg R$ abbiamo che il campo gravitazionale è debole e segue che
\begin{align}
	g_{00} = -1 - \frac{2 \Phi}{c^2},
\end{align}
pertanto, ci aspettiamo nella regione di buffer che l'espansione in equazione \eqref{eq:expansion_buffer} sia ancora valida, quindi si ha che la componente temporale del tensore metrico è data da
\begin{align}
	g_{00}(\vec{x}) = -1 + \frac{2M}{r} + \frac{3}{r^3} Q^{ij} n^{\langle i} n^{j \rangle} + o \left( \frac{R^3}{r^3} \right) - E_{ij} r^2 n^{\langle i} n^{j \rangle} + o \left( \frac{r^3}{\mathcal{R}^3} \right) + o \left( \frac{M^2}{r^2} \right),
\end{align}
in cui si noti che stiamo anche trascurando i termini di alto grado dell'espansione di campo debole, che sono dell'ordine di $\frac{M^2}{r^2}$. In questo caso, inoltre, segue che
\begin{align}
	E_{ij} = -\frac{1}{2} \frac{\partial^2 g_{00}}{\partial x^i \partial x^j} = R_{0i0j},
\end{align}
in quanto, 

	\chapter{Conclusioni}

In questo elaborato è stata illustrata l'efficacia dell'espansione in multipoli nella gravitazione newtoniana per descrivere in modo rigoroso le deviazioni dalla simmetria sferica, confrontando la decomposizione in armoniche sferiche con il formalismo cartesiano dei tensori STF. L'introduzione dei momenti di multipolo ha fornito degli strumenti analitici per quantificare la deviazione dei corpi celesti dalla simmetria sferica e, applicato ai dati gravimetrici terrestri, questo formalismo permette di quantificare lo schiacciamento rotazionale del pianeta, sperimentalmente dominato dal momento di quadrupolo, e di isolare, nei termini di ordine superiore, le anomalie locali generate dalla complessa geodinamica interna. È stato inoltre approfondito il problema delle deformazioni mareali indotte da sorgenti esterne, dimostrando, sia in ambito classico sia nel limite relativistico, come il campo mareale induca un momento di quadrupolo proporzionale al campo stesso, parametrizzato tramite i numeri di Love gravitazionali, e abbiamo mostrato il peculiare annullamento esatto del coefficiente di deformabilità mareale nel caso dei buchi neri.


I principali punti di forza di questo impianto teorico risiedono nella sua semplicità analitica e nell'ampia generalità applicativa, tuttavia emergono precisi vincoli fisici e operativi: dal punto di vista della modellizzazione, la determinazione dei coefficienti richiede la conoscenza dettagliata della distribuzione di densità e, nel contesto mareale, dell'equazione di stato della materia interna, necessaria per integrare le equazioni di campo all'interno del corpo ed eseguire il corretto raccordo sulla superficie. Inoltre, l'espansione in multipoli mostra evidenti limiti di convergenza quando applicata a configurazioni che si discostano marcatamente dalla simmetria sferica, come le galassie a spirale: in tali scenari, sebbene lo sviluppo resti formalmente valido, la lenta convergenza della serie impone l'inclusione di ordini molto elevati, rendendo il calcolo analitico poco pratico e richiedendo delle necessarie strategie algoritmiche per eliminare le instabilità numeriche connesse al calcolo delle funzioni associate di Legendre.

	\appendix
	\chapter{Calcolo Ricorsivo delle Funzioni di Legendre Associate ad alto ordine}

Nella geodesia, le espansioni troncate in armoniche sferiche di una funzione (come il potenziale gravitazionale) o delle sue derivate si riducono al calcolo di somme totali $S^{(d)}$. La notazione $d=0$ indica la funzione non derivata, mentre $d=1$ indica la derivata prima rispetto alla distanza polare $\theta$. La somma $S^{(d)}$ è definita come la somma dei contributi associati a ciascun ordine $m$:
\begin{equation}
S^{(d)} = c \sum_{m=0}^{M} \Omega_m^{(d)} \label{eq:somma_totale}
\end{equation}
dove $c$ è una costante e $M$ è il grado massimo dell'espansione. Il termine $\Omega_m^{(d)}$ raccoglie al suo interno le dipendenze dalla longitudine e la sommatoria, per il grado $n$, del prodotto tra i coefficienti geopotenziali normalizzati $\bar{E}_{nm\alpha}$ e le Funzioni di Legendre Associate (ALF) completamente normalizzate $\bar{P}_{nm}^{(d)}(\theta)$.

Il problema computazionale emerge nel calcolo di $\bar{P}_{nm}(\theta)$ per modelli ad altissimo grado (es. $M=2700$). Avvicinandosi ai poli ($\theta \rightarrow 0$ oppure $\theta \rightarrow \pi$), il valore assoluto $|\bar{P}_{nm}(\theta)|$ copre un intervallo di circa 5000 ordini di grandezza. Tale range viola i limiti dell'aritmetica a doppia precisione IEEE 754 (che ammette valori tra $\sim 10^{-310}$ e $\sim 10^{310}$) e i valori calcolati subiscono \textit{underflow} (troncamento a zero) o \textit{overflow} (errore di limite superato), invalidando le componenti $\Omega_m^{(d)}$ e impedendo il calcolo di $S^{(d)}$. Spiegheremo, dunque, rifacendosi a \cite{Holmes2002}, l'origine di questa instabilità numerica e come sia possibile recuperare la stabilità numerica ad alti valori di $l$.

\section{Origine dell'Instabilità nella Ricorsione Standard e risoluzione}

I valori di $\bar{P}_{nm}(\theta)$ vengono calcolati tramite equazioni ricorsive: il metodo più comune è la ricorsione "in avanti per colonna" (\textit{standard forward column recursion}), che incrementa il grado $n$ mantenendo costante l'ordine $m$
\begin{equation}
\bar{P}_{nm}(\theta) = a_{nm} t \bar{P}_{n-1,m}(\theta) - b_{nm} \bar{P}_{n-2,m}(\theta), \quad \forall n > m \label{eq:standard_col}
\end{equation}
dove $t = \cos \theta$, mentre $a_{nm}$ e $b_{nm}$ sono coefficienti indipendenti dalla latitudine.

La ricorsione necessita di "semi" iniziali: i valori settoriali dove $n=m$ e sono definiti come:
\begin{equation}
\bar{P}_{mm}(\theta) = u^m \Pi_m, \quad \forall m \ge 1
\end{equation}
dove $u = \sin \theta$ e $\Pi_m$ è un prodotto numericamente stabile. L'instabilità è dovuta esclusivamente al fattore $u^m$: in prossimità dei poli $u \rightarrow 0$ e per $m=2700$, il termine $u^m$ è infinitesimo e va istantaneamente in underflow, precludendo l'avvio della ricorsione.

Per aggirare l'underflow, l'algoritmo modificato rinuncia al calcolo diretto di $\bar{P}_{nm}(\theta)$ a favore della quantità normalizzata $\frac{\bar{P}_{nm}(\theta)}{u^m}$.
La validità algebrica di questa operazione si basa sull'architettura della ricorsione stessa: osservando l'equazione (\ref{eq:standard_col}), il calcolo itera esclusivamente sul grado $n$, mantenendo invariato l'ordine $m$ lungo l'intera "colonna". Poiché $m$ è costante durante questa fase, il termine $u^m = \sin^m \theta$ funge da mera costante moltiplicativa. 

Dividendo ambo i membri dell'Equazione (\ref{eq:standard_col}) per tale costante, le proprietà invariantive assicurano che la forma della legge ricorsiva non venga alterata:
\begin{equation}
\frac{\bar{P}_{nm}(\theta)}{u^m} = a_{nm} t \frac{\bar{P}_{n-1,m}(\theta)}{u^m} - b_{nm} \frac{\bar{P}_{n-2,m}(\theta)}{u^m}, \quad \forall n > m \label{eq:mod_col}
\end{equation}
In questo modo, la macchina itera su grandezze da cui è stato tolto il termine di decadimento esponenziale, il quale garantisce la totale stabilità senza distorcere la validità dell'algoritmo originario. 

\section{Scala Globale e la Reintroduzione tramite Schema di Horner}

La rimozione di $u^m$ salva i calcoli dall'underflow, ma spalanca le porte al rischio opposto (overflow) nei gradi estremi: infatti, per trovare il valore dell'armonica, siamo costretti a moltiplicare per $u^m$ nuovamente, il quale può essere gigantesco. Per arginare questo problema, le funzioni iniziali vengono moltiplicate per una costante di scala globale pari a $10^{-280}$. Essendo la ricorsione (\ref{eq:mod_col}) lineare, questo scalamento si trasmette a tutti i termini senza alterarne le proporzioni.

Una volta calcolate tutte le funzioni scalate, esse vengono moltiplicate per i coefficienti $\bar{E}_{nm\alpha}$ e sommate su $n$ per ottenere le componenti scalate dell'ordine $m$, che definiamo formalmente come $\frac{\Omega_m^{(d)}}{u^m}$ (sempre decurtate del fattore $10^{-280}$).

A questo punto è necessario calcolare la somma totale $S^{(d)}$ dell'equazione (\ref{eq:somma_totale}). Come già anticipato, occorrerebbe calcolare $\sum_{m=0}^{M} \left( \frac{\Omega_m^{(d)}}{u^m} \cdot u^m \right)$, ma il calcolo esplicito di $u^m$ per $m$ elevati genererebbe un overflow numrico: per evitare ciò, possiamo implementare lo \emph{Schema di Horner}. Tale algoritmo valuta i polinomi riducendo il calcolo delle potenze a una serie di moltiplicazioni a cascata; esplicitamente possiamo riscrivere la somma totale in base alla variabile $u$, l'espressione diventa:
\begin{equation}
S^{(d)} = c \left[ \left\{ \dots \left( \left[ \left\{ \frac{\Omega_M^{(d)}}{u^M} \right\} u + \frac{\Omega_{M-1}^{(d)}}{u^{M-1}} \right] u + \dots \right) u + \frac{\Omega_1^{(d)}}{u^1} \right\} u + c\Omega_0^{(d)} \right] \label{eq:horner}
\end{equation}
Questo algoritmo ripristina la variabile $u$ ``progressivamente'': partendo dal termine di grado massimo $M$, la somma viene accumulata e moltiplicata per $u$ \textit{step by step} ad ogni passaggio verso il grado inferiore. La dipendenza trigonometrica viene così reintegrata un "gradino" alla volta, garantendo che le somme parziali risiedano stabilmente nei limiti della doppia precisione. Al termine del ciclo, è sufficiente moltiplicare il risultato per $10^{280}$ (invertendo il fattore di scala iniziale) per giungere al valore esatto della somma $S^{(d)}$.


	\phantomsection
	\addcontentsline{toc}{chapter}{Bibliografia}
	\bibliographystyle{unsrt}
	\bibliography{bibliography.bib}
	\addcontentsline{toc}{chapter}{Ringraziamenti}
	\chapter{Ringraziamenti}

Questo rappresenta per me il capitolo più difficile da scrivere: rendere a parole e trasmettere un concetto su un qualcosa su cui studi ogni giorno, sbattendoci la testa fino ad vedere quel barlume di illuminazione che ti fa capire effettivamente cosa sta succedendo, è molto più facile, almeno a mio avviso, piuttosto che mettere su carta quello che vorresti dire a tutte quelle persone che, silenziosamente o meno, ti hanno accompagnato in questo lungo percorso. \\
Vorrei ringraziare, innanzitutto, il mio relatore, il professore Leonardo Gualtieri, per la sua disponibilità ad essere il mio relatore, oltre ad aver rappresentato, almeno per me, un docente fantastico, di cui ho avuto l'immensa fortuna di averlo sia come esercitatore del corso di Fisica 1 e sia come docente titolare del corso di Meccanica Quantistica, che ritengo essere il corso che ho apprezzato di più, sia per argomenti trattati e sia per la trattazione, nel percorso triennale di Fisica. \\
Subito dopo, vorrei ringraziare la mia famiglia, che mi ha supportato costantemente in questi anni, senza mai dubitare delle mie capacità: ringrazio i miei genitori, Antonietta (perdonami mamma se ti chiamo col nome intero) ed Enrico, per tutto il supporto morale che mi avete dato in questi anni; in particolar modo la mamma, che fra le mie numerose sfide e le sue numerose sfide e responsabilità, è sempre stata lì, nei momenti di bisogno, ad ascoltarmi nei (numerosi) momenti di sconforto e disperazione senza mai farmi pesare nulla. Successivamente, non posso non ringraziare la zia Adele e lo zio Maurizio, che mi hanno sempre supportato quasi come se fossi figlio loro, e mia cugina Giada, che ha sempre creduto nel suo cuginetto (anche se, purtroppo, i ``compiti'' di matematica che faccio adesso non puoi più correggerli facilmente come facevi un tempo). Infine, ringrazio la nonna Enrica, a cui ho realizzato il grande sogno di far vedere entrambi i nipoti laureati, e ringrazio sempre per la grande mano che ha dato a me e mia mamma in questi tre anni. \\
Non sarei riuscito a farcela senza le grandi amicizie che ho fatto in questi anni, di cui, ovviamente, devo ringraziare i ``quarratini'' OG, ovvero Carmine, Dario e Ale, che ormai conosco da tempo immemore, che hanno avuto modo di sperimentare la mia grandissima pazzia già prima di Fisica e hanno visto il suo andamento esponenziale in questi anni di ``cattività pisana'' in ogni chiamata su Discord. Ringrazio (fra virgolette anche se spero non me ne abbiate a male) i quarratini ``aggiunti'', Matteo, Sam, Michael e Tommaso, qua presenti, che ho conosciuto dopo, ma non di meno (in particolar modo Matteo che a Ferragosto di ritorno da Lavagna mi ha salvato la vita) che sto piano piano corrompendo verso la lega dei campioni \\
Ringrazio i miei coinquilini Anna e Matteo, con cui ho iniziato questo percorso $3$ anni fa, anche se il nostro ``cammino'' insieme è iniziato ben prima, quando sono arrivato il primo giorno di terza superiore al liceo scientifico. Adesso siamo tutti separati ahimé, ma, nonostante le normali discussioni che ci possono essere fra coinquilini, mi avete sempre sostenuto nei numerosi momenti del bisogno, spesso dinanzi ad un bel pasto dato da una ``minoranza'' (nuovo nome ufficiale di Just Eat).  \\
E ora iniziano i ringraziamenti dei ''pisani``: ringrazio Ismaele Bartoluccio, il mio daddy topologico, che è certamente una fra le persone più ``ganze'' che io abbia mai conosciuto (alla fine sono convenuto per questo termine, l'unico che può forse rendere su carta la grandezza dello spirito libero e indipendente di Isma): in questi anni mi ha preso sotto la sua ala, spingendomi nello studio (disperatissimo) delle aree della matematica più affascinanti ed esotiche per un fisico; oltre, nell'ultimo anno, ad avermi dato un supporto morale per tutta una serie di vicende della mia vita che, in questa sede, è meglio non raccontare, ma che Ismaele sa bene. Ringrazio poi Riccardo, che mi ha supportato e sopportato ogni giorno, ogni ora e ogni minuto di questa triennale, subendo i miei grandissimi danni psicologici ``alla 6pek'' dovuti alla mia pazzia (forse saranno stati questi che ti hanno convinto a fare Teorica insieme, piuttosto che Astrofisica ahah). Ringrazio Alessandro Terranova, che rimarrà il mio solo e unico presidentessimo, per avermi sostenuto più e più volte nei momenti in cui ero più vulnerabile, oltre a farmi conoscere il meraviglioso mondo di AI-SF (non Associazione Italiana Studio del Fegato obv) che, ad entrambi, rappresenta una grande famiglia estesa (sia per le grandissimi emozioni che ci procura e sia per i grandi momenti di tribolazione che ci ha procurato). Ringrazio il mitico Francesco Zeno Costanzo (o più brevemente il Gen) per avermi preso, a suon di insulti per le interlinee assenti ``alla Bolo'', sotto la sua ala pythonica e per avermi fatto capire quanto sia importante scrivere una stringa di spiegazione sia fondamentale (anche per la più banale delle funzioni) e per intarmi la \emph{toplane} (caro si scherza). Ringrazio Alberto Otera (il mio goat) per avermi fatto anche lui da psicologo dopo le lezioni pazze ed incredibili di Gruppi (mettersi accanto forse, come dicevi te, non era forse una buona idea) e di MQA, oltre a graziarmi con le sue mosse divine su League of Legends. Ringrazio Francesco Angelo Fabiano Antonacci per aver fatto parte in questo grandissimo esperimento sociale che consisteva nel porre, in un gruppo di laboratorio, la persona più teorica del dipartimento con, forse, la persona più sperimentale del dipartimento e per non aver mai perso la sua grandissima capacità nel dare i nomi nei file utili a laboratorio (ricordo a tutti l'esplicativo \emph{four\_int\_der\_train2.py}) e ringrazio Benedetta Santoro  

	\end{document}
